% Options for packages loaded elsewhere
\PassOptionsToPackage{unicode}{hyperref}
\PassOptionsToPackage{hyphens}{url}
\PassOptionsToPackage{dvipsnames,svgnames,x11names}{xcolor}
\documentclass[
  10pt,
]{article}
\usepackage{xcolor}
\usepackage[margin=2.5cm]{geometry}
\usepackage{amsmath,amssymb}
\setcounter{secnumdepth}{5}
\usepackage{iftex}
\ifPDFTeX
  \usepackage[T1]{fontenc}
  \usepackage[utf8]{inputenc}
  \usepackage{textcomp} % provide euro and other symbols
\else % if luatex or xetex
  \usepackage{unicode-math} % this also loads fontspec
  \defaultfontfeatures{Scale=MatchLowercase}
  \defaultfontfeatures[\rmfamily]{Ligatures=TeX,Scale=1}
\fi
\usepackage{lmodern}
\ifPDFTeX\else
  % xetex/luatex font selection
  \setmainfont[]{Arial Unicode MS}
  \setmonofont[]{Menlo}
\fi
% Use upquote if available, for straight quotes in verbatim environments
\IfFileExists{upquote.sty}{\usepackage{upquote}}{}
\IfFileExists{microtype.sty}{% use microtype if available
  \usepackage[]{microtype}
  \UseMicrotypeSet[protrusion]{basicmath} % disable protrusion for tt fonts
}{}
\makeatletter
\@ifundefined{KOMAClassName}{% if non-KOMA class
  \IfFileExists{parskip.sty}{%
    \usepackage{parskip}
  }{% else
    \setlength{\parindent}{0pt}
    \setlength{\parskip}{6pt plus 2pt minus 1pt}}
}{% if KOMA class
  \KOMAoptions{parskip=half}}
\makeatother
\usepackage{color}
\usepackage{fancyvrb}
\newcommand{\VerbBar}{|}
\newcommand{\VERB}{\Verb[commandchars=\\\{\}]}
\DefineVerbatimEnvironment{Highlighting}{Verbatim}{commandchars=\\\{\}}
% Add ',fontsize=\small' for more characters per line
\newenvironment{Shaded}{}{}
\newcommand{\AlertTok}[1]{\textcolor[rgb]{1.00,0.00,0.00}{\textbf{#1}}}
\newcommand{\AnnotationTok}[1]{\textcolor[rgb]{0.38,0.63,0.69}{\textbf{\textit{#1}}}}
\newcommand{\AttributeTok}[1]{\textcolor[rgb]{0.49,0.56,0.16}{#1}}
\newcommand{\BaseNTok}[1]{\textcolor[rgb]{0.25,0.63,0.44}{#1}}
\newcommand{\BuiltInTok}[1]{\textcolor[rgb]{0.00,0.50,0.00}{#1}}
\newcommand{\CharTok}[1]{\textcolor[rgb]{0.25,0.44,0.63}{#1}}
\newcommand{\CommentTok}[1]{\textcolor[rgb]{0.38,0.63,0.69}{\textit{#1}}}
\newcommand{\CommentVarTok}[1]{\textcolor[rgb]{0.38,0.63,0.69}{\textbf{\textit{#1}}}}
\newcommand{\ConstantTok}[1]{\textcolor[rgb]{0.53,0.00,0.00}{#1}}
\newcommand{\ControlFlowTok}[1]{\textcolor[rgb]{0.00,0.44,0.13}{\textbf{#1}}}
\newcommand{\DataTypeTok}[1]{\textcolor[rgb]{0.56,0.13,0.00}{#1}}
\newcommand{\DecValTok}[1]{\textcolor[rgb]{0.25,0.63,0.44}{#1}}
\newcommand{\DocumentationTok}[1]{\textcolor[rgb]{0.73,0.13,0.13}{\textit{#1}}}
\newcommand{\ErrorTok}[1]{\textcolor[rgb]{1.00,0.00,0.00}{\textbf{#1}}}
\newcommand{\ExtensionTok}[1]{#1}
\newcommand{\FloatTok}[1]{\textcolor[rgb]{0.25,0.63,0.44}{#1}}
\newcommand{\FunctionTok}[1]{\textcolor[rgb]{0.02,0.16,0.49}{#1}}
\newcommand{\ImportTok}[1]{\textcolor[rgb]{0.00,0.50,0.00}{\textbf{#1}}}
\newcommand{\InformationTok}[1]{\textcolor[rgb]{0.38,0.63,0.69}{\textbf{\textit{#1}}}}
\newcommand{\KeywordTok}[1]{\textcolor[rgb]{0.00,0.44,0.13}{\textbf{#1}}}
\newcommand{\NormalTok}[1]{#1}
\newcommand{\OperatorTok}[1]{\textcolor[rgb]{0.40,0.40,0.40}{#1}}
\newcommand{\OtherTok}[1]{\textcolor[rgb]{0.00,0.44,0.13}{#1}}
\newcommand{\PreprocessorTok}[1]{\textcolor[rgb]{0.74,0.48,0.00}{#1}}
\newcommand{\RegionMarkerTok}[1]{#1}
\newcommand{\SpecialCharTok}[1]{\textcolor[rgb]{0.25,0.44,0.63}{#1}}
\newcommand{\SpecialStringTok}[1]{\textcolor[rgb]{0.73,0.40,0.53}{#1}}
\newcommand{\StringTok}[1]{\textcolor[rgb]{0.25,0.44,0.63}{#1}}
\newcommand{\VariableTok}[1]{\textcolor[rgb]{0.10,0.09,0.49}{#1}}
\newcommand{\VerbatimStringTok}[1]{\textcolor[rgb]{0.25,0.44,0.63}{#1}}
\newcommand{\WarningTok}[1]{\textcolor[rgb]{0.38,0.63,0.69}{\textbf{\textit{#1}}}}
\setlength{\emergencystretch}{3em} % prevent overfull lines
\providecommand{\tightlist}{%
  \setlength{\itemsep}{0pt}\setlength{\parskip}{0pt}}
\usepackage{fvextra}
\DefineVerbatimEnvironment{Highlighting}{Verbatim}{breaklines,breakanywhere,commandchars=\\\{\}}
\DefineVerbatimEnvironment{verbatim}{Verbatim}{breaklines,breakanywhere}
\usepackage{bookmark}
\IfFileExists{xurl.sty}{\usepackage{xurl}}{} % add URL line breaks if available
\urlstyle{same}
\hypersetup{
  pdftitle={Tutorial: from zero to a running beacon},
  colorlinks=true,
  linkcolor={Maroon},
  filecolor={Maroon},
  citecolor={Blue},
  urlcolor={Blue},
  pdfcreator={LaTeX via pandoc}}

\title{Tutorial: from zero to a running beacon}
\author{}
\date{}

\begin{document}
\maketitle

{
\setcounter{tocdepth}{3}
\tableofcontents
}
\section{Tutorial: from zero to a running
beacon}\label{tutorial-from-zero-to-a-running-beacon}

Start with nothing. Finish with a GA4GH Beacon v2 service running on
your machine, holding data, answering genomic queries correctly.

\textbf{Time:} about 15 minutes, most of it waiting for a Docker build.

Every command and every output on this page was run end to end on
2026-08-19. Where output is shown, that is what was actually printed ---
not what should be printed. If yours differs, something is wrong and the
troubleshooting section at the end probably covers it.

\begin{quote}
\textbf{Deploying to a server instead?} This tutorial covers a local
instance. For a public deployment on a VM with TLS and a tunnel, see PR
\#9 (\texttt{docs/DEPLOYMENT\_TUTORIAL.md}), written from a real
production deploy. It is not merged yet, so read it from the pull
request.
\end{quote}

\subsection{What you will have at the
end}\label{what-you-will-have-at-the-end}

\begin{itemize}
\tightlist
\item
  MongoDB 5, Redis 6 and the Django API running in Docker
\item
  100 variants, 50 individuals and 1 dataset loaded
\item
  A beacon that answers \texttt{exists:\ true} for a variant it holds
  and \texttt{false} for one it does not
\item
  Enough understanding of the query path to change it safely
\end{itemize}

\subsection{Before you start}\label{before-you-start}

You need \textbf{Docker} with Compose v2, and about 2 GB of free disk.

\begin{Shaded}
\begin{Highlighting}[]
\ExtensionTok{docker} \AttributeTok{{-}{-}version}
\ExtensionTok{docker}\NormalTok{ info }\OperatorTok{\textgreater{}}\NormalTok{/dev/null }\KeywordTok{\&\&} \BuiltInTok{echo} \StringTok{"daemon is running"}
\end{Highlighting}
\end{Shaded}

Verified against Docker 29.2.1. You do \textbf{not} need Python, MongoDB
or Node installed locally --- everything runs in containers. You will
want Python later to run the test suites, but not for this tutorial.

\begin{center}\rule{0.5\linewidth}{0.5pt}\end{center}

\subsection{Step 1 --- Get the code}\label{step-1-get-the-code}

\begin{Shaded}
\begin{Highlighting}[]
\FunctionTok{git}\NormalTok{ clone git@github.com:AfriGen{-}D/variant{-}checker{-}beacon.git}
\BuiltInTok{cd}\NormalTok{ afrigen{-}beacon{-}v2}
\end{Highlighting}
\end{Shaded}

\textbf{Clone somewhere Docker Desktop can share} --- under your home
directory is safe. On macOS, \texttt{/tmp} is \emph{not} in Docker
Desktop's shared paths by default, and bind mounts from there silently
resolve to empty directories inside the VM instead of failing. That
produces a container that cannot write its log file and an API that
never starts, with an error that points nowhere near the cause. This
cost an hour to diagnose while writing this page.

\subsection{Step 2 --- Check the environment file
exists}\label{step-2-check-the-environment-file-exists}

The dev stack reads \texttt{.env.boolean} from the repository root.
\textbf{It is gitignored (\texttt{.gitignore:62}), so a fresh clone does
not have it} and \texttt{docker\ compose} will refuse to start:

\begin{Shaded}
\begin{Highlighting}[]
\NormalTok{env file /path/to/repo/.env.boolean not found}
\end{Highlighting}
\end{Shaded}

Create it from the example:

\begin{Shaded}
\begin{Highlighting}[]
\FunctionTok{cp}\NormalTok{ .env.example .env.boolean}
\end{Highlighting}
\end{Shaded}

\textbf{Then make one edit, or nothing will work.} \texttt{.env.example}
is a \emph{production} template: it sets
\texttt{SECURE\_SSL\_REDIRECT=True}, which makes Django 301-redirect
every request to \texttt{https://localhost:8000} --- where nothing is
listening. You get a stack that looks healthy in \texttt{docker\ ps} and
answers nothing.

\begin{Shaded}
\begin{Highlighting}[]
\FunctionTok{sed} \AttributeTok{{-}i} \StringTok{\textquotesingle{}\textquotesingle{}} \StringTok{\textquotesingle{}s/\^{}SECURE\_SSL\_REDIRECT=True/SECURE\_SSL\_REDIRECT=False/\textquotesingle{}}\NormalTok{ .env.boolean   }\CommentTok{\# macOS}
\CommentTok{\# GNU/Linux: sed {-}i \textquotesingle{}s/\^{}SECURE\_SSL\_REDIRECT=True/SECURE\_SSL\_REDIRECT=False/\textquotesingle{} .env.boolean}
\end{Highlighting}
\end{Shaded}

\texttt{DJANGO\_SECRET\_KEY} ships as \texttt{CHANGE\_ME}, which is fine
locally. Never reuse it anywhere real.

\begin{quote}
If you change \texttt{.env.boolean} later,
\texttt{docker\ compose\ restart} will \textbf{not} pick it up ---
\texttt{env\_file} is read when the container is created. Use
\texttt{docker\ compose\ -f\ compose/docker-compose.dev.yml\ up\ -d\ -\/-force-recreate\ beacon-api}.
\end{quote}

\subsection{Step 3 --- Start the stack}\label{step-3-start-the-stack}

\begin{Shaded}
\begin{Highlighting}[]
\ExtensionTok{docker}\NormalTok{ compose }\AttributeTok{{-}f}\NormalTok{ compose/docker{-}compose.dev.yml up }\AttributeTok{{-}d} \AttributeTok{{-}{-}build}\NormalTok{ mongodb redis beacon{-}api}
\end{Highlighting}
\end{Shaded}

The first run builds the API image and takes a few minutes. Compose
starts MongoDB and Redis first and \textbf{waits for both to report
healthy} before starting the API, so you should not see connection
errors during startup.

Expected tail:

\begin{Shaded}
\begin{Highlighting}[]
\NormalTok{ Container beacon{-}dev{-}mongodb  Healthy}
\NormalTok{ Container beacon{-}dev{-}redis    Healthy}
\NormalTok{ Container beacon{-}api{-}boolean  Started}
\end{Highlighting}
\end{Shaded}

Confirm all three are up:

\begin{Shaded}
\begin{Highlighting}[]
\ExtensionTok{docker}\NormalTok{ compose }\AttributeTok{{-}f}\NormalTok{ compose/docker{-}compose.dev.yml ps}
\end{Highlighting}
\end{Shaded}

\begin{Shaded}
\begin{Highlighting}[]
\NormalTok{NAME                 IMAGE                STATUS                    PORTS}
\NormalTok{beacon{-}api{-}boolean   compose{-}beacon{-}api   Up (health: starting)     127.0.0.1:8000{-}\textgreater{}8000/tcp}
\NormalTok{beacon{-}dev{-}mongodb   mongo:5.0            Up (healthy)              127.0.0.1:27017{-}\textgreater{}27017/tcp}
\NormalTok{beacon{-}dev{-}redis     redis:6{-}alpine       Up (healthy)              127.0.0.1:6380{-}\textgreater{}6379/tcp}
\end{Highlighting}
\end{Shaded}

Note the ports bind to \texttt{127.0.0.1}, not \texttt{0.0.0.0} --- the
stack is not reachable from your network. Redis is on \textbf{6380}
locally to avoid colliding with any Redis you already run.

\subsubsection{If Step 3 fails with ``address pools have been fully
subnetted''}\label{if-step-3-fails-with-address-pools-have-been-fully-subnetted}

\begin{Shaded}
\begin{Highlighting}[]
\NormalTok{failed to create network walk\_beacon{-}dev{-}network: Error response from daemon:}
\NormalTok{all predefined address pools have been fully subnetted}
\end{Highlighting}
\end{Shaded}

Docker has run out of private subnets, not disk or memory. Each compose
project claims one and they are not released until the network is
removed. On a machine running several projects you will hit this before
you hit any resource limit.

Find an idle network belonging to something you own and remove it:

\begin{Shaded}
\begin{Highlighting}[]
\ExtensionTok{docker}\NormalTok{ network ls}
\ExtensionTok{docker}\NormalTok{ network inspect }\OperatorTok{\textless{}}\NormalTok{name}\OperatorTok{\textgreater{}}\NormalTok{ {-}{-}format }\StringTok{\textquotesingle{}\{\{.Name\}\}: \{\{len .Containers\}\} attached\textquotesingle{}}
\ExtensionTok{docker}\NormalTok{ network rm }\OperatorTok{\textless{}}\NormalTok{name}\OperatorTok{\textgreater{}}          \CommentTok{\# only if 0 attached, and only if it is yours}
\end{Highlighting}
\end{Shaded}

Do \textbf{not} reach for \texttt{docker\ network\ prune}. It removes
every unused network on the machine, including ones other people's
stopped stacks will want back.

\subsection{Step 4 --- Confirm the API is
alive}\label{step-4-confirm-the-api-is-alive}

Give it about ten seconds after the container starts, then:

\begin{Shaded}
\begin{Highlighting}[]
\ExtensionTok{curl} \AttributeTok{{-}s}\NormalTok{ http://localhost:8000/api/health}
\end{Highlighting}
\end{Shaded}

\begin{Shaded}
\begin{Highlighting}[]
\FunctionTok{\{}\DataTypeTok{"status"}\FunctionTok{:}\StringTok{"healthy"}\FunctionTok{,}\DataTypeTok{"version"}\FunctionTok{:}\StringTok{"2.0.0{-}boolean"}\FunctionTok{,}
 \DataTypeTok{"services"}\FunctionTok{:\{}\DataTypeTok{"database"}\FunctionTok{:}\StringTok{"healthy"}\FunctionTok{,}\DataTypeTok{"cache"}\FunctionTok{:}\StringTok{"healthy"}\FunctionTok{\}\}}
\end{Highlighting}
\end{Shaded}

Both \texttt{database} and \texttt{cache} must say \texttt{healthy}. If
either does not, the API started before its dependencies were ready ---
\texttt{docker\ compose\ restart\ beacon-api} and try again.

\begin{quote}
A passing health check proves the API can reach MongoDB and Redis. It
proves \textbf{nothing} about whether queries return correct answers.
That distinction has caused real incidents on this project;
\texttt{docs/DEPLOYMENT\_PREREQUISITES.md} explains why at length.
\end{quote}

\subsection{Step 5 --- Query the empty
beacon}\label{step-5-query-the-empty-beacon}

Before loading anything, ask it something. This matters: you want to see
what ``no'' looks like from a beacon that genuinely holds nothing, so
you can tell it apart from a broken one later.

\begin{Shaded}
\begin{Highlighting}[]
\ExtensionTok{curl} \AttributeTok{{-}s} \StringTok{\textquotesingle{}http://localhost:8000/api/g\_variants?assemblyId=GRCh38\&referenceName=1\&start=1000000\textquotesingle{}}
\end{Highlighting}
\end{Shaded}

\begin{Shaded}
\begin{Highlighting}[]
\FunctionTok{\{}\DataTypeTok{"meta"}\FunctionTok{:} \FunctionTok{\{}\ErrorTok{...}\FunctionTok{\},}
 \DataTypeTok{"responseSummary"}\FunctionTok{:} \FunctionTok{\{}\DataTypeTok{"exists"}\FunctionTok{:} \KeywordTok{false}\FunctionTok{,} \DataTypeTok{"numTotalResults"}\FunctionTok{:} \DecValTok{0}\FunctionTok{\},}
 \DataTypeTok{"response"}\FunctionTok{:} \FunctionTok{\{}\DataTypeTok{"resultSets"}\FunctionTok{:} \OtherTok{[]}\FunctionTok{,} \DataTypeTok{"beaconHandovers"}\FunctionTok{:} \OtherTok{[}\ErrorTok{...}\OtherTok{]}\FunctionTok{\}\}}
\end{Highlighting}
\end{Shaded}

\textbf{Read that response shape carefully.} The answer is at
\texttt{responseSummary.exists}. There is no top-level \texttt{exists}
field, and several older documents in this repository show client
examples that parse one. Those examples silently report ``not found''
for everything. See \url{api-contract.md}.

\subsection{Step 6 --- Load the test
data}\label{step-6-load-the-test-data}

\begin{Shaded}
\begin{Highlighting}[]
\ExtensionTok{docker}\NormalTok{ compose }\AttributeTok{{-}f}\NormalTok{ compose/docker{-}compose.dev.yml exec }\AttributeTok{{-}T}\NormalTok{ beacon{-}api }\DataTypeTok{\textbackslash{}}
\NormalTok{  python manage.py load\_boolean\_test\_data }\AttributeTok{{-}{-}settings}\OperatorTok{=}\NormalTok{beacon\_project.settings\_boolean}
\end{Highlighting}
\end{Shaded}

\begin{Shaded}
\begin{Highlighting}[]
\NormalTok{  {-} Individuals: 50}

\NormalTok{You can now test Boolean queries!}
\end{Highlighting}
\end{Shaded}

Check what landed:

\begin{Shaded}
\begin{Highlighting}[]
\ExtensionTok{docker}\NormalTok{ compose }\AttributeTok{{-}f}\NormalTok{ compose/docker{-}compose.dev.yml exec }\AttributeTok{{-}T}\NormalTok{ mongodb }\DataTypeTok{\textbackslash{}}
\NormalTok{  mongo beacon\_db }\AttributeTok{{-}{-}quiet} \AttributeTok{{-}{-}eval} \StringTok{\textquotesingle{}db.getCollectionNames().forEach(function(c)\{ print("  "+c+": "+db[c].count()) \})\textquotesingle{}}
\end{Highlighting}
\end{Shaded}

\begin{Shaded}
\begin{Highlighting}[]
\NormalTok{  datasets: 1}
\NormalTok{  individuals: 50}
\NormalTok{  query\_logs: 62}
\NormalTok{  variants: 100}
\end{Highlighting}
\end{Shaded}

The \texttt{mongo:5.0} image ships the \textbf{legacy \texttt{mongo}
shell}. \texttt{mongosh} does not exist in it and any command using
\texttt{mongosh} will fail --- some older docs in this repo get this
wrong.

\texttt{query\_logs} is the audit trail; it fills up as you query.
Client IPs in it are truncated to a /24 before being written, and rows
expire on a TTL index. See \texttt{beacon\_api/privacy.py}.

\subsection{Step 7 --- Find a variant that actually
exists}\label{step-7-find-a-variant-that-actually-exists}

Do not guess coordinates. Ask the database what it holds:

\begin{Shaded}
\begin{Highlighting}[]
\ExtensionTok{docker}\NormalTok{ compose }\AttributeTok{{-}f}\NormalTok{ compose/docker{-}compose.dev.yml exec }\AttributeTok{{-}T}\NormalTok{ mongodb }\DataTypeTok{\textbackslash{}}
\NormalTok{  mongo beacon\_db }\AttributeTok{{-}{-}quiet} \AttributeTok{{-}{-}eval} \StringTok{\textquotesingle{}var v=db.variants.findOne();}
\StringTok{  print("assembly="+v.assembly\_id+"  chr="+v.reference\_name+"  start="+v.start+"  ref="+v.reference\_bases+"  alt="+v.alternate\_bases)\textquotesingle{}}
\end{Highlighting}
\end{Shaded}

\begin{Shaded}
\begin{Highlighting}[]
\NormalTok{assembly=GRCh38  chr=chr1  start=42497823  ref=A  alt=G}
\end{Highlighting}
\end{Shaded}

Yours will differ --- use your own values below.

Notice the chromosome is stored as \texttt{chr1}, not \texttt{1}. Which
form ends up in the database depends on the ingest pipeline that wrote
it, which is exactly why the query path matches both.

\subsection{Step 8 --- Ask the beacon a real
question}\label{step-8-ask-the-beacon-a-real-question}

\begin{Shaded}
\begin{Highlighting}[]
\ExtensionTok{curl} \AttributeTok{{-}s} \StringTok{\textquotesingle{}http://localhost:8000/api/g\_variants?assemblyId=GRCh38\&referenceName=1\&start=42497823\textquotesingle{}} \DataTypeTok{\textbackslash{}}
  \KeywordTok{|} \ExtensionTok{python3} \AttributeTok{{-}c} \StringTok{\textquotesingle{}import json,sys; print(json.load(sys.stdin)["responseSummary"])\textquotesingle{}}
\end{Highlighting}
\end{Shaded}

\begin{Shaded}
\begin{Highlighting}[]
\NormalTok{\{\textquotesingle{}exists\textquotesingle{}: True, \textquotesingle{}numTotalResults\textquotesingle{}: 1\}}
\end{Highlighting}
\end{Shaded}

\textbf{That is a working beacon.} You asked with a bare \texttt{1} and
it matched data stored as \texttt{chr1}.

Now confirm it says no when it should:

\begin{Shaded}
\begin{Highlighting}[]
\ExtensionTok{curl} \AttributeTok{{-}s} \StringTok{\textquotesingle{}http://localhost:8000/api/g\_variants?assemblyId=GRCh38\&referenceName=1\&start=999\textquotesingle{}} \DataTypeTok{\textbackslash{}}
  \KeywordTok{|} \ExtensionTok{python3} \AttributeTok{{-}c} \StringTok{\textquotesingle{}import json,sys; print(json.load(sys.stdin)["responseSummary"])\textquotesingle{}}
\end{Highlighting}
\end{Shaded}

\begin{Shaded}
\begin{Highlighting}[]
\NormalTok{\{\textquotesingle{}exists\textquotesingle{}: False, \textquotesingle{}numTotalResults\textquotesingle{}: 0\}}
\end{Highlighting}
\end{Shaded}

A negative control is not optional here. A beacon that answers
\texttt{true} to everything looks identical to a working one until
someone checks.

\subsection{Step 9 --- See the vocabulary
handling}\label{step-9-see-the-vocabulary-handling}

\texttt{hg38} and \texttt{GRCh38} are the \textbf{same genome build} in
different vocabularies --- UCSC and GRC. A beacon must answer both
identically:

\begin{Shaded}
\begin{Highlighting}[]
\ControlFlowTok{for}\NormalTok{ a }\KeywordTok{in}\NormalTok{ GRCh38 hg38 HG38}\KeywordTok{;} \ControlFlowTok{do}
  \BuiltInTok{printf} \StringTok{\textquotesingle{}\%{-}8s \textquotesingle{}} \StringTok{"}\VariableTok{$a}\StringTok{"}
  \ExtensionTok{curl} \AttributeTok{{-}s} \StringTok{"http://localhost:8000/api/g\_variants?assemblyId=}\VariableTok{$a}\StringTok{\&referenceName=1\&start=42497823"} \DataTypeTok{\textbackslash{}}
    \KeywordTok{|} \ExtensionTok{python3} \AttributeTok{{-}c} \StringTok{\textquotesingle{}import json,sys; print(json.load(sys.stdin)["responseSummary"]["exists"])\textquotesingle{}}
\ControlFlowTok{done}
\end{Highlighting}
\end{Shaded}

\begin{Shaded}
\begin{Highlighting}[]
\NormalTok{GRCh38   True}
\NormalTok{hg38     True}
\NormalTok{HG38     True}
\end{Highlighting}
\end{Shaded}

And an assembly the beacon cannot answer for is refused rather than
answered:

\begin{Shaded}
\begin{Highlighting}[]
\ExtensionTok{curl} \AttributeTok{{-}s} \StringTok{\textquotesingle{}http://localhost:8000/api/g\_variants?assemblyId=GRCh99\&referenceName=1\&start=42497823\textquotesingle{}}
\end{Highlighting}
\end{Shaded}

\begin{Shaded}
\begin{Highlighting}[]
\FunctionTok{\{}\DataTypeTok{"error"}\FunctionTok{:} \FunctionTok{\{}\DataTypeTok{"errorCode"}\FunctionTok{:} \DecValTok{400}\FunctionTok{,}
  \DataTypeTok{"errorMessage"}\FunctionTok{:} \StringTok{"Unknown assembly: GRCh99. Recognised builds: GRCh37, GRCh38. This beacon may not hold data for all of them."}\FunctionTok{\}\}}
\end{Highlighting}
\end{Shaded}

\subsubsection{Why this step exists}\label{why-this-step-exists}

Until 2026-08-19 the middle line read \texttt{hg38\ \ \ False}. The
query path compared the caller's literal spelling against what was
stored, so a caller using UCSC vocabulary got a confident ``no'' for a
variant the panel was holding. Nobody reported it, because nothing
looked broken --- a researcher simply concluded the variant was absent
from African reference data. It was two clicks away in the UI, because
GRCh37 was then one of only two options in the assembly dropdown.

Both halves of that are now closed. \texttt{hg38} canonicalises to
\texttt{GRCh38}, and a build this beacon holds no data for is
\textbf{refused rather than answered}:

\begin{Shaded}
\begin{Highlighting}[]
\ExtensionTok{curl} \AttributeTok{{-}s} \StringTok{\textquotesingle{}http://localhost:8000/api/g\_variants?assemblyId=GRCh37\&referenceName=1\&start=42497823\textquotesingle{}}
\end{Highlighting}
\end{Shaded}

\begin{Shaded}
\begin{Highlighting}[]
\FunctionTok{\{}\DataTypeTok{"error"}\FunctionTok{:} \FunctionTok{\{}\DataTypeTok{"errorCode"}\FunctionTok{:} \DecValTok{501}\FunctionTok{,}
  \DataTypeTok{"errorMessage"}\FunctionTok{:} \StringTok{"This beacon holds no data for assembly GRCh37, so it cannot}
\StringTok{   answer for that build. This is not a statement about whether the variant}
\StringTok{   exists. Data held: GRCh38."}\FunctionTok{\}\}}
\end{Highlighting}
\end{Shaded}

A 501 is a signal a client can act on; \texttt{exists:\ false} is an
answer it will believe.

Note what the 400 does \textbf{not} say. It lists the builds this beacon
\emph{recognises}, and then explicitly declines to promise it holds data
for them. An earlier wording said ``this beacon answers for GRCh37,
GRCh38'' --- which sent a caller who mistyped an assembly straight into
the 501 above, because GRCh37 is recognised and unheld.
\texttt{assembly.py} has no access to the dataset catalogue, so rather
than claim coverage it cannot verify, it stops claiming.

The assembly dropdown now offers only GRCh38, so a UI user cannot reach
this --- it is for API callers.

That is the failure mode this project guards hardest against, and the
rule it produced is the one to carry into your own changes: \textbf{when
the beacon cannot answer the question you asked, it must refuse ---
never answer a different question and return 200.}
\texttt{beacon\_api/filters.py} is the reference implementation of that
judgement and its docstring explains the reasoning;
\texttt{beacon\_api/assembly.py} is the fix that came out of it.

\subsubsection{One more thing to notice}\label{one-more-thing-to-notice}

Look closely at a positive answer:

\begin{Shaded}
\begin{Highlighting}[]
\NormalTok{\{\textquotesingle{}exists\textquotesingle{}: True, \textquotesingle{}numTotalResults\textquotesingle{}: 0\}}
\end{Highlighting}
\end{Shaded}

\texttt{exists} is \texttt{True} while \texttt{numTotalResults} is
\texttt{0}. That is not a typo --- the counter reports matching
\emph{datasets} under some paths rather than matching variants. Do not
build a client that infers absence from \texttt{numTotalResults}.

\subsection{Step 9b --- Ask the same question by
POST}\label{step-9b-ask-the-same-question-by-post}

Beacon v2's own request format nests the parameters. Every client that
follows the spec sends this shape, so it is worth seeing it work:

\begin{Shaded}
\begin{Highlighting}[]
\ExtensionTok{curl} \AttributeTok{{-}s} \AttributeTok{{-}X}\NormalTok{ POST http://localhost:8000/api/g\_variants }\DataTypeTok{\textbackslash{}}
  \AttributeTok{{-}H} \StringTok{\textquotesingle{}Content{-}Type: application/json\textquotesingle{}} \DataTypeTok{\textbackslash{}}
  \AttributeTok{{-}d} \StringTok{\textquotesingle{}\{"query":\{"requestParameters":\{"assemblyId":"GRCh38","referenceName":"1","start":\textless{}POS\textgreater{}\}\}\}\textquotesingle{}} \DataTypeTok{\textbackslash{}}
  \KeywordTok{|} \ExtensionTok{python3} \AttributeTok{{-}c} \StringTok{\textquotesingle{}import json,sys; print(json.load(sys.stdin)["responseSummary"])\textquotesingle{}}
\end{Highlighting}
\end{Shaded}

Substitute the \texttt{\textless{}POS\textgreater{}} you found in Step
7. It must return exactly what the GET in Step 8 returned --- if POST
and GET ever disagree, one of them is wrong.

\subsubsection{Why this step exists, and what it nearly
shipped}\label{why-this-step-exists-and-what-it-nearly-shipped}

Until 2026-08-20 this returned
\texttt{400\ Invalid\ characters\ detected\ in\ query}. The sanitizer
stringified the nested dict before pattern-matching, and the Python
repr's own quote matched its injection pattern --- an injection verdict
on punctuation the sanitizer introduced itself. GET was unaffected,
because flat query strings never produce a repr, which is why nothing
looked broken.

Fixing that alone would have been worse than the bug. Nothing in the
codebase read \texttt{query.requestParameters}, so once the body was let
through the query carried no filters at all: an impossible locus
returned \texttt{exists:\ true} with \texttt{numTotalResults:\ 100},
identical to an empty body. An honest error would have become a
confident YES for a variant that cannot exist.

Both halves shipped together. The lesson is worth more than the fix:
\textbf{a change that unblocks a path is not safe until you check what
is behind the path.}

\subsection{Step 10 --- Look around the rest of the
API}\label{step-10-look-around-the-rest-of-the-api}

\begin{Shaded}
\begin{Highlighting}[]
\ExtensionTok{curl} \AttributeTok{{-}s}\NormalTok{ http://localhost:8000/api/          }\KeywordTok{|} \FunctionTok{head} \AttributeTok{{-}c}\NormalTok{ 200   }\CommentTok{\# beacon info}
\ExtensionTok{curl} \AttributeTok{{-}s}\NormalTok{ http://localhost:8000/api/datasets                  }\CommentTok{\# what is loaded}
\ExtensionTok{curl} \AttributeTok{{-}s}\NormalTok{ http://localhost:8000/api/entry\_types               }\CommentTok{\# supported entry types}
\ExtensionTok{curl} \AttributeTok{{-}s}\NormalTok{ http://localhost:8000/api/map                       }\CommentTok{\# endpoint map}
\ExtensionTok{curl} \AttributeTok{{-}s} \StringTok{\textquotesingle{}http://localhost:8000/api/individuals?sex=MALE\textquotesingle{}}
\end{Highlighting}
\end{Shaded}

\subsection{Step 11 --- Run the test
suites}\label{step-11-run-the-test-suites}

The query-correctness suites need \textbf{no dependencies at all} --- no
Django, no database, no network:

\begin{Shaded}
\begin{Highlighting}[]
\ExtensionTok{python3} \AttributeTok{{-}m}\NormalTok{ unittest }\DataTypeTok{\textbackslash{}}
\NormalTok{  beacon\_api.test\_query\_semantics }\DataTypeTok{\textbackslash{}}
\NormalTok{  beacon\_api.test\_query\_injection }\DataTypeTok{\textbackslash{}}
\NormalTok{  beacon\_api.test\_pagination\_filters }\DataTypeTok{\textbackslash{}}
\NormalTok{  beacon\_api.test\_assembly }\DataTypeTok{\textbackslash{}}
\NormalTok{  beacon\_api.test\_query\_vocabulary }\DataTypeTok{\textbackslash{}}
\NormalTok{  beacon\_api.test\_capabilities }\DataTypeTok{\textbackslash{}}
\NormalTok{  beacon\_api.test\_release }\DataTypeTok{\textbackslash{}}
\NormalTok{  beacon\_api.test\_request\_body}
\end{Highlighting}
\end{Shaded}

\begin{Shaded}
\begin{Highlighting}[]
\NormalTok{Ran 200 tests in 0.002s}

\NormalTok{OK}
\end{Highlighting}
\end{Shaded}

That is the merge gate. If you change how a query is interpreted, these
are the tests that must stay green --- and the ones you extend.

\texttt{beacon\_api.test\_middleware} is different: it imports Django,
so it needs a Python between \textbf{3.9 and 3.12}. Django 4.0 cannot
import on 3.13+, which removed the \texttt{cgi} module it depends on.

\subsection{Step 12 --- Shut down}\label{step-12-shut-down}

\begin{Shaded}
\begin{Highlighting}[]
\ExtensionTok{docker}\NormalTok{ compose }\AttributeTok{{-}f}\NormalTok{ compose/docker{-}compose.dev.yml down}
\end{Highlighting}
\end{Shaded}

Add \texttt{-v} to delete the MongoDB volume as well and start
completely fresh next time.

\begin{center}\rule{0.5\linewidth}{0.5pt}\end{center}

\subsection{Troubleshooting}\label{troubleshooting}

\textbf{\texttt{database:\ unhealthy} in the health response.} The API
started before MongoDB finished initialising.
\texttt{docker\ compose\ -f\ compose/docker-compose.dev.yml\ restart\ beacon-api}.

\textbf{Port 8000, 27017 or 6380 already allocated.} Something else is
using it. Find it with \texttt{lsof\ -i\ :8000}, or change the host side
of the mapping in \texttt{compose/docker-compose.dev.yml}.

\textbf{\texttt{mongosh:\ not\ found}.} The \texttt{mongo:5.0} image
only has the legacy \texttt{mongo} shell. Use \texttt{mongo}.

\textbf{Every query returns \texttt{exists:\ false} after loading data.}
Almost always a coordinate or vocabulary mismatch rather than a broken
beacon. Check three things, in order: are you using the assembly the
data was loaded under (Step 7 prints it); is your \texttt{start} the
stored 0-based value rather than the 1-based VCF position; and does the
variant you are asking about actually exist (ask Mongo directly, as in
Step 7).

\textbf{A query changed answers after you edited code.} Redis caches
responses for five minutes.
\texttt{docker\ compose\ -f\ compose/docker-compose.dev.yml\ exec\ redis\ redis-cli\ FLUSHDB}.

\textbf{Everything returns HTTP 301 / \texttt{curl} shows
\texttt{Location:\ https://localhost:8000}.}
\texttt{SECURE\_SSL\_REDIRECT=True} is still set in
\texttt{.env.boolean}. See Step 2 --- and remember the container must be
recreated, not restarted, for an env change to apply.

\textbf{\texttt{ValueError:\ Unable\ to\ configure\ handler\ \textquotesingle{}file\textquotesingle{}}
in the API logs, container never becomes healthy.} The container runs as
the non-root user \texttt{beacon} (uid 1000,
\texttt{Dockerfile.boolean:39}) and cannot write
\texttt{logs/beacon.log}. Almost always means the repository is
somewhere Docker Desktop does not share --- see Step 1. If you are on a
shared path, Docker creates \texttt{logs/} with workable ownership
automatically and you should not hit this.

\textbf{Two clones of this repo fight over the same containers.} The
Compose project name is derived from the directory holding the compose
file, which is \texttt{compose/} in every clone --- so two checkouts
both resolve to a project called \texttt{compose} and share containers
\emph{and} the MongoDB volume. Pass \texttt{-p\ beacon-myfeature} to
\texttt{docker\ compose} in a second checkout, or you will be querying
the other one's data without noticing.

\textbf{The build fails installing Python dependencies.} The backend is
pinned to Django 4.0.10 by \texttt{django-mongoengine}, and needs
\texttt{setuptools\textless{}70}. The Dockerfile handles the ordering;
if you are installing by hand outside Docker, install
\texttt{setuptools\textless{}70} first.

\subsection{Where to go next}\label{where-to-go-next}

\begin{itemize}
\tightlist
\item
  \url{architecture.md} --- how a query becomes an answer
\item
  \url{testing.md} --- what runs, and what gates a merge
\item
  \url{api-contract.md} --- the response envelope, and the trap in it
\item
  \url{contributing.md} --- branching, commits, PRs
\end{itemize}

Loading \textbf{real} genomic data rather than the test fixture is a
separate job: see \texttt{docs/ILIFU\_DATA\_LOADING\_GUIDE.md} and
\texttt{afrigend-beacon2-tools/README.md}. Read the data-tools caveats
in \url{testing.md} before running an import --- there are known issues
with loading a second dataset over a first.

\end{document}
